% ============================================================
%  v3_conclusion.tex — Section 7: Conclusion
%  Soccer-only version
% ============================================================

\section{Conclusion}
\label{sec:conclusion}

FIFA's 2022 stoppage-time directive provides a rare natural
experiment in time regulation: a global regulator mandated fuller
enforcement of an existing rule, and the world watched the
demonstration live. We trace the directive's effects through five
layers of the game.

\textit{Compliance was near-universal.} Between 95\% and 100\% of
referees in treated leagues increased their stoppage time, depending
on the minimum-match threshold applied. The directive achieved this
without formal sanctions, fines, or monitoring beyond what already
existed. The mechanism was focal-point coordination: a public signal
that shifted the common expectation of acceptable enforcement
practice, producing convergence in referee behaviour as measured by
a declining coefficient of variation.

\textit{The first stage is large and robust.} Treated leagues added
0.74~minutes of second-half stoppage relative to controls ($p =
0.007$; wild cluster bootstrap $p = 0.036$). The effect is stable
across three DiD specifications, confirmed by Callaway-Sant'Anna
estimation, and present in every treated league except Serie~A.

\textit{Goals were displaced, not created.} Total goals per match
were statistically unchanged (DiD $\beta = -0.033$, $p = 0.39$).
Scoring in the post-90th-minute window rose by 20\%, while the
76--84 minute window declined by 8.6\%. Per-minute scoring rates
were constant across all time bins. The directive purchased
additional minutes of play; those minutes produced goals at exactly
the historical rate.

\textit{The game's dramatic arc shifted.} The share of matches in
which the result changed during stoppage time rose from 7.3\% to
10.3\%---approximately 55~additional decisive moments per Big~Five
season. Late equalizers rose by 23\%, late winners by 17\%.
Competitive balance, as measured by within-league points dispersion
and normalised Herfindahl indices, was unaffected.

\textit{The quantity-intensity decomposition is clean.} Across
goals, fouls, and cards, the dominant mechanism is a pure quantity
effect: more minutes at unchanged per-minute rates. The only
behavioural signal is a declining per-minute foul rate, which the
DiD attributes to a secular trend rather than a directive-specific
response.

\medskip

Two limitations deserve emphasis. First, the directive's adoption
was simultaneous across all five treated leagues, not staggered.
Our identification rests on the treated-versus-never-treated
contrast, not on variation in adoption timing. Second, the first
stage is modest in absolute terms---45~seconds per match---because
control leagues also increased stoppage, likely through shared
refereeing bodies. The DiD estimate is a lower bound on the
directive's full effect.

These findings speak to a question that extends well beyond football:
what happens when a regulator credibly enforces a rule that was
previously under-enforced? In this setting, the answer is
straightforward. Agents complied near-universally, output expanded
proportionally along the extensive margin, and per-unit productivity
was unchanged. The simplest model---more time, same rate---explains
the data. Whether this linearity holds at larger doses, or in
settings where fatigue accumulates faster, remains an open question
for future work.
